\documentclass[aspectratio=169,10pt]{beamer}

\usetheme{Madrid}
\usecolortheme{seahorse}
\setbeamertemplate{navigation symbols}{}
\setbeamertemplate{footline}[frame number]

\usepackage{amsmath,amssymb,booktabs,array}
\usepackage{tikz}
\usetikzlibrary{arrows.meta,positioning,shapes.geometric,fit,calc}
\usepackage{xcolor}
\usepackage{hyperref}

\definecolor{darkblue}{RGB}{0,65,120}
\hypersetup{colorlinks=true,linkcolor=darkblue,urlcolor=darkblue}

\newcommand{\code}[1]{\texttt{#1}}
\newcommand{\TODO}[1]{\textcolor{red}{\textbf{TODO:}} \emph{#1}}

\title[L07: PINN]{L07 -- Physics-informed Neural Networks (PINN)}
\subtitle{Foundations Lecture Series}
\author{Course Instructor}
\institute{Microgrid 3-phase Security AI Project}
\date{Foundations Lecture 7}

\begin{document}

\begin{frame}
  \titlepage
  \begin{center}
    \small Keywords: physics-informed loss, soft constraints, limit-aware learning, Raissi 2019
  \end{center}
\end{frame}

% ============================================================
\begin{frame}{Lecture roadmap}
\begin{enumerate}
  \item The PINN idea -- physics as a soft constraint
  \item Loss design: data + physics
  \item Strict vs soft constraints
  \item Limit-aware learning -- the project's flavor
  \item When PINNs help and when they don't
  \item Connection to the project's Week 6 PI-MLP
\end{enumerate}
\vspace{0.6em}
\begin{block}{Prerequisite}
L05 (MLPs, loss design, PyTorch). Helpful: L01--L04 for the physics being constrained.
\end{block}
\end{frame}

% ============================================================
\section*{1. The PINN idea}

\begin{frame}{Where PINNs come from}
\TODO{Briefly trace the lineage from prior physics-based ML to Raissi-Perdikaris-Karniadakis 2019. State the central idea: encode the known PDE/ODE/algebraic constraint into the loss function so the model is penalized for violating physics, not just for fitting data poorly.}
\end{frame}

\begin{frame}{Why bother with physics in the loss}
\TODO{Three reasons: (1) extrapolation beyond the training distribution, (2) data efficiency, (3) interpretability. Concrete example from microgrid security: a model that respects voltage limits ``for free'' is more trustworthy than one that learns them from data alone.}
\end{frame}

% ============================================================
\section*{2. Loss design}

\begin{frame}{Data loss + physics loss}
\TODO{Write $\mathcal{L} = \mathcal{L}_{\text{data}} + \lambda \mathcal{L}_{\text{phys}}$. Explain that $\mathcal{L}_{\text{phys}}$ is typically the squared residual of a physical equation evaluated at the model's prediction. Discuss the role and sensitivity of $\lambda$.}
\end{frame}

\begin{frame}{Multi-task framing}
\TODO{Many PINNs in power systems are really multi-task: classification + regression + ranking + limit penalty. Show the project's Week 6 PI-MLP composite loss as a concrete example.}
\end{frame}

\begin{frame}{Choosing physics terms}
\TODO{Concrete physics terms relevant to this project: voltage limit penalty, current limit penalty, VUF penalty, power-balance residual, sequence-component consistency. Note that not all of these are strict PINN-style PDE residuals -- some are just engineering limits.}
\end{frame}

% ============================================================
\section*{3. Strict vs soft constraints}

\begin{frame}{Soft constraints (the usual PINN)}
\TODO{Add a penalty to the loss. Model can still violate the constraint; the penalty just makes that costly. Easy to implement, gradient-friendly. The default approach.}
\end{frame}

\begin{frame}{Strict constraints (projection / hard layer)}
\TODO{Project the output onto the feasible set, or build a network architecture that cannot violate the constraint by construction. Harder to design, often necessary for safety-critical applications.}
\end{frame}

\begin{frame}{Which one for screening?}
\TODO{For pre-screening (the project's framing), soft constraints are usually fine -- we still re-run the simulator on flagged contingencies. For deployment without simulator backup, you would want hard constraints or conformal prediction.}
\end{frame}

% ============================================================
\section*{4. Limit-aware learning}

\begin{frame}{The project's flavor of physics-informed}
\TODO{The project does \emph{not} encode the full AC power-flow residual into the loss. It uses ``limit-aware'' soft penalties: voltage limits, current limits, VUF threshold. This is a pragmatic subset of full PINN. Be honest with students about this.}
\end{frame}

\begin{frame}{Severity regression as physics-informed}
\TODO{Predicting not just label (safe / unsafe) but also severity is a form of physics-informed multi-task learning. The severity head learns continuous physical quantities, which can stabilize the classification head.}
\end{frame}

\begin{frame}{Ranking loss as physics-informed}
\TODO{The project also adds a pairwise ranking loss: for top-k screening, the order of risk scores matters more than absolute probabilities. Frame this as another physics-aware multi-task component.}
\end{frame}

% ============================================================
\section*{5. When PINNs help (and don't)}

\begin{frame}{When PINNs help}
\TODO{Sparse data, known physics, extrapolation, safety-critical applications. Mention published cases from fluid mechanics, biomechanics, and power systems.}
\end{frame}

\begin{frame}{When PINNs do not help (or hurt)}
\TODO{Lots of data + well-tuned baseline. Wrong physics (e.g., a constraint that conflicts with the actual data-generating process). Overweighted penalty term. Poor optimization landscape from competing loss terms.}
\end{frame}

% ============================================================
\section*{6. Connection to the project}

\begin{frame}{What you will see in Week 6}
\TODO{Show the high-level PI-MLP loss decomposition: BCE + severity MSE + component-margin + consistency + ranking. Promise to walk through each term. Forward-reference the ablation table from the paper.}
\end{frame}

% ============================================================
\begin{frame}{Homework}
\TODO{Suggest: take the MLP from L05 (e.g., trained on a small N-1 dataset), add a voltage-limit penalty term, retune $\lambda$, and report whether recall and precision improve. Plot loss components over training.}
\end{frame}

\begin{frame}{Exit checklist}
\TODO{Items: (1) write the generic PINN loss form, (2) distinguish soft and strict constraints, (3) list the 5 physics terms used in the project, (4) name one scenario where PINNs help and one where they don't.}
\end{frame}

\begin{frame}{References}
\TODO{Raissi, Perdikaris, Karniadakis 2019 (\emph{JCP}); Karniadakis et al. 2021 \emph{Nature Reviews Physics} survey; Krishnapriyan et al. 2021 (\emph{NeurIPS}, on PINN failure modes); Recent power-system PINN surveys (Misyris et al., Stiasny et al.).}
\end{frame}

\end{document}
